\documentclass[10pt,a4paper]{amsart}
\usepackage[margin=19mm]{geometry}
\usepackage{amsmath,amssymb,mathtools}
\usepackage[scr=rsfs,cal=euler]{mathalfa}
\usepackage{xfrac}
\usepackage[unicode,hidelinks,bookmarks=false]{hyperref}
\newcommand{\ie}{i.e.\ }
\newcommand{\cf}{cf.\ }
\newcommand{\resp}{respectively\ }
\newcommand{\Subsection}{\subsection}
\providecommand{\nobreakdash}{\nobreak\hbox{-}}
\setlength{\emergencystretch}{2em}
\multlinegap0pt
\usepackage[normalem]{ulem}
\usepackage{enumerate}
\usepackage{csquotes}
\usepackage{amssymb}
\usepackage{mathtools}
\usepackage{xcolor}
\usepackage{tikz}
\newtheorem{theorem}{Theorem}
\newtheorem{lemma}{Lemma}
\title{Enumerating monophyletic characters in mathematical phylogenetics}
\author{Mareike Fischer, Michael Hendriksen, Kristina Wicke}
\date{Source: arXiv:2607.22097v1 (CC BY 4.0)}
\begin{document}
\maketitle

\noindent\emph{Source and licence.} Enumerating monophyletic characters in mathematical phylogenetics. Version arXiv:2607.22097v1 (CC BY 4.0), distributed by arXiv under the Creative Commons Attribution 4.0 International licence, http://creativecommons.org/licenses/by/4.0/, as recorded in the arXiv licence metadata for this exact version and shown on the abstract page https://arxiv.org/abs/2607.22097v1. Full licence notice: \texttt{licenses/arxiv-2607.22097-CC-BY-4.0.txt}.\par\smallskip

\noindent\textit{Editorial scope.} Counted unit: the article's theorem thm:binary.summary (source0.tex lines 916--937). The article's complete original proof of it is delivered with it (source0.tex lines 940--985, one \texttt{proof} environment of 46 lines). No mathematics is written, repaired or re-derived: the statement and the proof are the article's own lines, in the article's own notation, and no retained line is substituted or re-flowed.  Retained uncounted as the source-local context the proof needs: lines 457--467; lines 904--906.  Nothing was omitted from any retained span, so the package carries no omission record.  No retained line contains a citation, so the wrapper carries no bibliography and no bibliography entry.  No retained line contains a figure environment, an image-inclusion command or drawing code, so no image is delivered and the package declares no asset dependency.  Editorial formatting only, in addition: an AMS \texttt{amsart} preamble that restates the article's own declarations and macros used by the retained lines, namely the environments \texttt{theorem}, \texttt{lemma} on separate counters, together with the standard packages those lines need.  The retained environments print in this document as Theorem 1, Lemma 1 in order of appearance, whereas the article prints the same items with the numbering of its own full document.  The complete native source of the article as distributed in the arXiv e-print for this exact version is bundled unchanged as \texttt{sources/205911/source0.tex}.
\begin{theorem}
\label{thm:gen.theorem}
    Let $T$ be a phylogenetic tree on $n$ leaves and $C$ a set of traits with $|C|=c$. Let $\mathcal{A}$ be the set of antichains that exist in $T$. Then

    \[m_N(T,c) = \sum_{A \in \mathcal{A}} (-1)^{|A|+1} (c)_{|A|} (c-|A|)^{n-s(A)},\]
with the convention that $(c-|A|)^{n-s(A)}=1$ if both $c-|A|=0$ and $n-s(A)=0$. If, instead, $\mathcal{A}$ is the set of relevant antichains, then 

\[m(T,c) = \sum_{A \in \mathcal{A}} (-1)^{|A|+1} (c)_{|A|} (c-|A|)^{n-s(A)},\]

with the same convention.
\end{theorem}

\begin{lemma}\label{lem_len2anti}
    Let $T$ be a rooted binary phylogenetic tree with $n\geq 2$ leaves. Then there is a unique antichain of size $2$ that covers all the leaves of $T$, consisting of the two child vertices of the root of $T$.
\end{lemma}

\begin{theorem}\label{thm:binary.summary}
Let \(T=(T_1,T_2)\) be a rooted binary phylogenetic tree with
\(n\ge 2\) leaves, where \(T_1,T_2\) have \(n_1,n_2\) leaves respectively and \(n_1\ge n_2\). Then:

\begin{align*}
m(T,2) &=
\begin{cases}
2n-4,& n_2>1,\\
2n-2,& n_2=1,
\end{cases} \\
%%%
m_N(T,2)&=4n-4,\\
%%%
fm(T,2) &=
\begin{cases}
4,& n_2>1,\\
2,& n_2=1,
\end{cases} \qquad \text{and}\\
%%%
fm_N(T,2)&=4.
\end{align*}
\end{theorem}

\begin{proof}
By Theorem \ref{thm:gen.theorem} and the convention $(x)_k=0$ for $x<k$,
only antichains of size at most $2$ contribute when $c=2$.

Moreover, if $|A|=2$, then the term
\[
(c-|A|)^{n-s(A)}
\]
is non-zero only when $s(A)=n$, that is, when the antichain covers all leaves of $T$.

By Lemma \ref{lem_len2anti}, there is a unique antichain of size $2$ that covers all leaves, namely the set consisting of the two children of the root.

Now, each type of monophyly count differs only by which antichains are admissible

\begin{itemize}
\item non-relevant monophyly: all antichains;
\item relevant monophyly: relevant antichains, i.e., antichains that contain no leaves;
\item full monophyly: antichains that cover all leaves;
\item relevant full monophyly: relevant antichains that cover all leaves.
\end{itemize}

We first consider non-relevant monophyly. For non-relevant monophyly, every vertex of $T$ may form an antichain of size $1$. Since a rooted binary tree with $n$ leaves has $2n-1$ vertices, these antichains contribute
\[
2(2n-1)=4n-2.
\]
The unique covering antichain of size $2$ contributes
\[
-(2)_2= -2,
\]
and hence
\[
m_N(T,2)=4n-4.\]

We now consider non-relevant full monophyly. In this case, we must restrict the size-$1$ antichains to also specifically be those that cover all leaves of the tree, which means the only permitted size-$1$ antichain consists of the root. There are two possible trait assignments for descendants of the root, which together with the two possible assignments for the antichain of size $2$ from Lemma \ref{lem_len2anti} gives us that $fm_N(T,2)=4$.

For relevant monophyly, only internal vertices may appear in antichains. Since $T$ has $n-1$ internal vertices, antichains of size $1$ contribute
\[
2(n-1)=2n-2.
\] 

Again, by Lemma \ref{lem_len2anti}, there is only one possible size-$2$ antichain. However, if $n_2=1$, one of the vertices in this antichain will be a leaf, so this is not a permissible antichain. Hence if $n_2=1$, $m(T,2)=2n-2$. However, if $n_2>1$, it is permissible, and so by Theorem \ref{thm:gen.theorem}, we must subtract the two possible trait assignments, and in this case $m(T,2)=2n-4$.

Finally, we consider relevant full monophyly. In this case we are restricted to only antichains that cover the leaves of $T$, and furthermore do not contain leaves. There are therefore two possible antichains -- the size-$1$ antichain containing only the root, and the size-$2$ antichain consisting of the two children of the root. If $n_2=1$, then the latter contains a leaf, and hence the only valid relevant fully monophyletic characters are those that assign all descendants of the root the same character, of which there are $2$. Hence if $n_2=1$, we obtain $fm(T,2)=2$. If, however, $n_2>1$, the size-$2$ antichain is valid, with a further two assignments, and thus $fm(T,2)=4$.

This proves the theorem.
\end{proof}

\end{document}
